% !TeX program = pdflatex
% El pentágono, verificado por máquina — la relación de cinco términos de Abel en Lean 4.
% Secuela de "El reloj que nunca suena": cierra el hilo del five-term que aquel paper dejó abierto.
% Build: pdflatex five-term-lean-es.tex (dos veces, para referencias)
% Verificado contra el fuente Lean 2026-06-13 (lake build Dilog.FiveTerm, exit 0;
% #print axioms en cada teorema -> propext, Classical.choice, Quot.sound).
% Toolchain v4.30.0-rc2, Mathlib pin 701fb6e9.
% Todas las instancias numéricas reverificadas a 30 dps (research/_tmp).

\documentclass[11pt]{article}
\usepackage[T1]{fontenc}
\usepackage[spanish,es-noshorthands]{babel}
\usepackage{lmodern}
\DeclareUnicodeCharacter{2082}{\ensuremath{_2}}
\usepackage{amsmath,amssymb,amsthm,booktabs,graphicx,hyperref}
\usepackage{xurl}
\usepackage[table]{xcolor}
\definecolor{ggteal}{HTML}{1B6F8C}
\definecolor{ggband}{HTML}{DCE7EB}
\definecolor{ggamber}{HTML}{E08A1E}
\usepackage[margin=1in]{geometry}
\usepackage{microtype}
\usepackage{float}
\usepackage[section]{placeins}
\usepackage[most]{tcolorbox}
\usepackage{tikz}
\usetikzlibrary{arrows.meta,positioning}
\setlength{\emergencystretch}{2em}
\newtheorem{theorem}{Teorema}
\newtheorem{lemma}{Lema}
\newtheorem{remark}{Observaci\'on}
\newcommand{\lean}[1]{\texttt{#1}}
\newcommand{\R}{\mathbb{R}}
\newcommand{\Li}{\mathrm{Li}_2}
\DeclareMathOperator{\Lrog}{L}
\hypersetup{colorlinks=true,linkcolor=ggteal,citecolor=ggteal,urlcolor=ggteal}
\newtcbox{\keybox}{on line,boxrule=0.9pt,colframe=ggteal,colback=ggband!40,
  arc=3pt,left=8pt,right=8pt,top=5pt,bottom=5pt}
\newcommand{\keyeq}[1]{\begin{center}\keybox{$\displaystyle #1$}\end{center}}
\newtcolorbox{recipebox}[1]{colframe=ggamber,colback=ggamber!8,arc=3pt,boxrule=0.8pt,
  title=\textbf{#1},coltitle=white,colbacktitle=ggamber}

\title{\bf El pent\'agono, verificado por m\'aquina\\[4pt]
  \large\normalfont La relaci\'on de cinco t\'erminos de Abel en Lean~4}
\author{Carles Mar\'in\thanks{Formalizado con asistencia de IA (Claude, Anthropic); la
matem\'atica y todas las afirmaciones son responsabilidad del autor. La metodolog\'ia y el
l\'imite exacto de la contribuci\'on est\'an en la Secci\'on~\ref{sec:boundary}.}\\[3pt]
  \normalsize Investigador independiente\\
  \normalsize \texttt{karlesmarin@gmail.com}}
\date{Junio 2026}

\begin{document}
\maketitle

\begin{abstract}
Presento una prueba verificada por m\'aquina, \lean{sorry}-libre, en Lean~4 sobre Mathlib,
de la \emph{relaci\'on de cinco t\'erminos de Abel} para el dilogaritmo. Con $\Lrog$ la
funci\'on $\Lrog$ de Rogers, dice, para $x,y\in(0,1)$,
\[
  \Lrog(x)+\Lrog(y)=\Lrog(xy)+\Lrog\!\Big(\tfrac{x(1-y)}{1-xy}\Big)
    +\Lrog\!\Big(\tfrac{y(1-x)}{1-xy}\Big).
\]
Esta sola identidad es toda la teor\'ia de peso~2 en una l\'inea: por teoremas de Wojtkowiak
y de~Jeu, \emph{toda} ecuaci\'on funcional de $\Li$ con argumentos racionales es
consecuencia suya; es la relaci\'on que define el grupo de Bloch, y el pent\'agono de las
\'algebras de cl\'uster y del dilogaritmo cu\'antico. Un paper compa\~nero, \emph{El reloj
que nunca suena}, construy\'o un dilogaritmo en Lean desde cero y subi\'o la escalera de la
raz\'on \'aurea \emph{sin} esta relaci\'on, dej\'andola expl\'icitamente abierta; este paper
cierra ese hilo. La prueba es el argumento cl\'asico de derivaci\'on hecho del todo riguroso:
fijando $y$, la diferencia de cinco t\'erminos tiene derivada id\'enticamente nula en
$(0,1)$ ---nueve t\'erminos logar\'itmicos colapsan a una base de cinco y el resto es
\lean{ring}--- luego es constante, y su valor en la frontera cuando $x\to0^+$ se calcula por
compresi\'on de $\log x\cdot\log(1-x)\to0$. Soy exacto sobre el l\'imite: la matem\'atica es
enteramente cl\'asica (Spence, Abel, Rogers; la universalidad de Wojtkowiak y de~Jeu); la
contribuci\'on es la formalizaci\'on ---hasta donde s\'e, la primera relaci\'on de cinco
t\'erminos verificada por m\'aquina, sobre el primer dilogaritmo verificado por m\'aquina, en
cualquier asistente de pruebas--- junto con tres bloques reutilizables. El teorema central y
sus bloques dependen solo de \lean{propext}, \lean{Classical.choice} y \lean{Quot.sound}.
\end{abstract}

\medskip
\noindent\emph{C\'omo leer este paper.} Es la historia de cuatro preguntas, cada una abriendo
una secci\'on y entregando al lector la siguiente, caminando de un enigma a un hilo cerrado.
El enigma: un paper compa\~nero alcanz\'o las evaluaciones m\'as profundas del dilogaritmo
\emph{sin} la identidad que todos llaman central ---as\'i que, \,¿merece siquiera la pena
verificarla por m\'aquina (Q1)? S\'i, porque no es una identidad m\'as sino el
\emph{generador} de todas; eso le gana la prueba (Q2), una prueba que la m\'aquina afil\'o
en un punto que la pluma pasa de puntillas (Q3); y una vez certificada, las piezas de gala
que el compa\~nero alcanz\'o por el camino largo caen como corolarios (Q4). Si solo lee una
cosa, lea el Teorema~\ref{thm:five} y la Figura~\ref{fig:pentagon}; el inventario Lean es la
Tabla~\ref{tab:inventory}; los l\'imites honestos est\'an en la Secci\'on~\ref{sec:boundary};
reproducir son tres comandos.

\begin{figure}[h]
\centering
\begin{tikzpicture}[
  obj/.style={draw=ggteal,thick,rounded corners=2pt,fill=ggband!40,align=center,
    font=\small,inner sep=5pt,minimum height=10mm},
  q/.style={font=\scriptsize\itshape,text=ggamber!80!black,align=center},
  arr/.style={-{Stealth[length=2.4mm]},thick,ggteal}]
  \node[obj] (clock) {el paper \emph{Reloj}\\sube sin Abel};
  \node[obj,right=24mm of clock] (gen) {el cinco-t\'erminos\\es el \emph{generador}};
  \node[obj,right=24mm of gen] (proof) {tres pasos:\\$G'\!=\!0$, const, frontera};
  \node[obj,below=15mm of proof] (caught) {lo que caz\'o\\la m\'aquina};
  \node[obj,left=24mm of caught] (coro) {escaleras \emph{con} Abel,\\Bloch, cl\'uster};
  \node[obj,left=24mm of coro,fill=ggamber!25,draw=ggamber] (done)
    {\textbf{el pent\'agono,}\\\textbf{verificado}};
  \draw[arr] (clock) -- node[q,above=1mm] {Q1: ¿vale\\la pena?} (gen);
  \draw[arr] (gen) -- node[q,above=1mm] {Q2: ¿misma\\receta humilde?} (proof);
  \draw[arr] (proof) -- node[q,right=1mm] {Q3: ¿qu\'e salt\'o\\la pluma?} (caught);
  \draw[arr] (caught) -- node[q,above=1mm] {Q4: ¿qu\'e cae\\gratis?} (coro);
  \draw[arr] (coro) -- (done);
\end{tikzpicture}
\caption{El mapa del lector. El hilo arranca donde el paper compa\~nero \emph{Reloj} lo dej\'o
(la escalera \'aurea subida sin Abel) y cierra en el pent\'agono verificado por m\'aquina;
cada flecha es la pregunta que responde su secci\'on. Cada nodo es un bloque de Lean
\lean{sorry}-libre.}
\label{fig:readermap}
\end{figure}

\section{Q1: ¿por qu\'e verificar por m\'aquina una identidad de la que el compa\~nero
prescindi\'o?}\label{sec:intro}

Hay identidades que importan por lo que afirman; esta importa por d\'onde no deja de
reaparecer. Tome los cinco n\'umeros que se sientan en las esquinas de un pent\'agono,
aplique a cada uno el dilogaritmo, y los cinco valores no son libres: obedecen una relaci\'on
fija, la misma cualquiera que sea el pent\'agono de partida. Ese solo hecho es, disfrazado,
la relaci\'on que define el grupo de Bloch en la geometr\'ia hiperb\'olica, el pent\'agono de
mutaciones en las \'algebras de cl\'uster, y la identidad del pent\'agono del dilogaritmo
cu\'antico en la f\'isica matem\'atica ---una misma relaci\'on con tres m\'ascaras. Tiene dos
siglos y ha sido redescubierta bajo una docena de nombres; y con todo, tras una d\'ecada de
uno de los esfuerzos de formalizaci\'on m\'as ambiciosos jam\'as emprendidos, a ning\'un
asistente de pruebas se le hab\'ia dicho nunca que es cierta. Este paper se lo dice ---y
luego pregunta por qu\'e algo as\'i est\'a en todas partes a la vez.

Un anticipo de lo que hace la relaci\'on, antes de toda teor\'ia. Ponga $x=y=\tfrac12$. Tras
$\tfrac{(1/2)(1/2)}{1-1/4}=\tfrac13$ la relaci\'on de abajo colapsa a
\[
  2\,\Lrog\!\big(\tfrac12\big)=\Lrog\!\big(\tfrac14\big)+2\,\Lrog\!\big(\tfrac13\big),
  \qquad\text{es decir}\qquad \Lrog\!\big(\tfrac14\big)+2\,\Lrog\!\big(\tfrac13\big)
  =\frac{\pi^2}{6},
\]
ya que $\Lrog(\tfrac12)=\pi^2/12$. Que una suma de tres valores del dilogaritmo en $\tfrac14$
y $\tfrac13$ sea \emph{exactamente} $\pi^2/6$ no es obvio; es una de las infinitas sombras de
una sola ley en dos variables ---y esa ley es lo que verificamos por m\'aquina aqu\'i.

\emph{Cuenta hist\'orica.} El dilogaritmo es de Euler (d\'ecada de 1760). Sus ecuaciones
funcionales las cartografiaron Landen (1780), Spence (1809), Abel (1827, publicado
en~1881~\cite{abel}) y Kummer; Rogers dio la forma sim\'etrica en 1907~\cite{rogers}. La
profundidad de la relaci\'on de cinco t\'erminos se volvi\'o estructural solo despu\'es, con
los reguladores de Bloch~\cite{bloch} y el grupo que lleva su nombre.

Hay un motivo para dudar. Un paper compa\~nero de esta l\'inea, \emph{El reloj que nunca
avanza}~\cite{clock}, construy\'o un dilogaritmo en Lean desde cero ---su serie, su derivada
$-\log(1-x)/x$, la reflexi\'on de Euler, la transformaci\'on de Landen, la f\'ormula de
duplicaci\'on--- y luego alcanz\'o las piezas de gala cl\'asicas del tema, las evaluaciones
de la raz\'on \'aurea $\Li(1/\varphi)$ y $\Li(1/\varphi^2)$, por un sistema lineal de tres
por tres, \emph{sin relaci\'on de cinco t\'erminos alguna}. Todo manual llega a esos valores
por Abel; la formalizaci\'on hall\'o una puerta lateral. As\'i que la primera pregunta
honesta no es c\'omo probar la relaci\'on de cinco t\'erminos sino si vale la pena: si las
piezas de gala se alcanzan sin ella, ¿en qu\'e se gana la identidad el sueldo?

La respuesta es que la relaci\'on de cinco t\'erminos no es una identidad entre varias. Es el
\emph{generador} de todo el tema. Wojtkowiak prob\'o que toda ecuaci\'on funcional de una
variable de $\Li$ con argumentos racionales es combinaci\'on $\mathbb{Z}$-lineal de
instancias de la relaci\'on de cinco t\'erminos; de~Jeu lo extendi\'o a cualquier n\'umero de
variables~\cite{wojtkowiak,dejeu}. La misma relaci\'on \emph{define} el grupo de Bloch, luego
gobierna los invariantes de congruencia por tijeras y los vol\'umenes de las $3$-variedades
hiperb\'olicas~\cite{bloch,suslin,neumannzagier}; en las \'algebras de cl\'uster es el
\emph{pent\'agono}, la relaci\'on de cinco mutaciones~\cite{nakanishi}; y en su forma
cu\'antica es la identidad del pent\'agono del dilogaritmo cu\'antico de
Faddeev--Kashaev~\cite{faddeevkashaev}. La puerta lateral del compa\~nero sube una escalera;
la relaci\'on de cinco t\'erminos es la escalinata entera. Eso es lo que merece certificarse.
\emph{¿Puede el mismo kit humilde que prob\'o la reflexi\'on y Landen alcanzar tambi\'en el
generador?}

El enunciado m\'as limpio usa la \emph{funci\'on $\Lrog$ de Rogers}
$\Lrog(x):=\Li(x)+\tfrac12\log x\log(1-x)$, que absorbe los logaritmos sueltos de la forma
cruda de $\Li$ y convierte la relaci\'on en una suma pura de cinco valores de $\Lrog$:
\keyeq{\Lrog(x)+\Lrog(y)=\Lrog(xy)+\Lrog\!\Big(\tfrac{x(1-y)}{1-xy}\Big)
  +\Lrog\!\Big(\tfrac{y(1-x)}{1-xy}\Big),\qquad x,y\in(0,1).}

\begin{figure}[t]
\centering
\begin{tikzpicture}[scale=1.5,every node/.style={font=\small}]
  \foreach \i/\lab in {0/{$x$},1/{$\tfrac{x(1-y)}{1-xy}$},2/{$y$},
      3/{$\tfrac{y(1-x)}{1-xy}$},4/{$xy$}}{
    \node (n\i) at ({90+\i*72}:1.25) {};
    \node at ({90+\i*72}:1.62) {\lab};
    \fill[ggteal] ({90+\i*72}:1.25) circle (1.5pt);
  }
  \foreach \i in {0,...,4}{
    \pgfmathtruncatemacro{\j}{mod(\i+1,5)}
    \draw[ggteal,-{Stealth[length=4pt]}] (n\i) -- (n\j);
  }
\end{tikzpicture}
\caption{\label{fig:pentagon}Los cinco argumentos sobre un ciclo. La simetr\'ia de orden
cinco de este pent\'agono es la sombra combinatoria que comparten la relaci\'on del grupo de
Bloch, el pent\'agono de mutaciones de las \'algebras de cl\'uster, y la identidad del
pent\'agono del dilogaritmo cu\'antico.}
\end{figure}

\section{Q2: ¿se alcanza el generador con la misma receta de tres pasos?}\label{sec:proof}

\emph{Cuenta hist\'orica.} La prueba por derivaci\'on es la m\'as antigua: fijar una
variable, mostrar que la derivada en la otra se anula, evaluar la constante en un extremo. Es
en esencia la de Abel, repetida en Lewin~\cite{lewin} y Zagier~\cite{zagier}. Toda la
cuesti\'on es si ``la derivada se anula'' y ``eval\'ua en el extremo'' sobreviven a hacerse
del todo formales.

\begin{theorem}[\lean{Dilog.rogersL\_fiveterm}]\label{thm:five}
Para todo $x,y\in(0,1)$,
\[
  \Lrog(x)+\Lrog(y)=\Lrog(xy)+\Lrog\!\Big(\tfrac{x(1-y)}{1-xy}\Big)
    +\Lrog\!\Big(\tfrac{y(1-x)}{1-xy}\Big).
\]
\end{theorem}

La receta es la que el paper compa\~nero us\'o para cada identidad elemental: derivar,
concluir que es constante, evaluar en la frontera. Fija $y\in(0,1)$ y pon
\[
  G(t)=\Lrog(t)-\Lrog(ty)-\Lrog\!\Big(\tfrac{t(1-y)}{1-ty}\Big)
    -\Lrog\!\Big(\tfrac{y(1-t)}{1-ty}\Big).
\]
El teorema es $G(x)=-\Lrog(y)$. Tres movimientos.

\paragraph{1. La derivada se anula.} Para $x,y\in(0,1)$ los tres argumentos compuestos
vuelven a estar en $(0,1)$, luego cada t\'ermino es derivable. Con $\Lrog'(t)=
-\tfrac12\big(\tfrac{\log(1-t)}{t}+\tfrac{\log t}{1-t}\big)$, la regla de la cadena escribe
$G'(t)$ como cuatro de estas expresiones pesadas por las derivadas de los argumentos
compuestos. Todo logaritmo que aparece ---de $t$, de $ty$, de los dos argumentos
fraccionarios y sus complementos--- se desarrolla, por $\log(ab)=\log a+\log b$ y
$\log(a/b)=\log a-\log b$, en la base de cinco elementos $\{\log t,\log y,\log(1-t),
\log(1-y),\log(1-ty)\}$. En esa base $G'(t)$ es una expresi\'on racional cuyos coeficientes
se cancelan todos: $G'\equiv0$. En Lean este es el \'unico paso delicado
(\lean{fiveterm\_hasDerivAt\_zero}): los nueve logaritmos se reescriben a la base y el resto
racional se limpia y se cierra con \lean{ring}.

\paragraph{2. Luego $G$ es constante.} Una funci\'on con derivada nula en el intervalo
abierto conexo $(0,1)$ es constante ah\'i ---el \lean{IsOpen.is\_const\_of\_deriv\_eq\_zero}
de Mathlib, el teorema del valor medio disfrazado. As\'i $G(x)=G(t)$ para todo $t\in(0,1)$.

\paragraph{3. El valor en la frontera.} Cuando $t\to0^+$, $ty\to0$,
$\tfrac{t(1-y)}{1-ty}\to0$ y $\tfrac{y(1-t)}{1-ty}\to y$, luego
$G(t)\to\Lrog(0)-\Lrog(0)-\Lrog(0)-\Lrog(y)=-\Lrog(y)$ ---\emph{siempre que} $\Lrog$ sea
continua por la derecha en $0$ con $\Lrog(0)=0$. Esta es la \'unica sutileza anal\'itica, la
misma que el paper compa\~nero aisl\'o. $\Lrog(t)=\Li(t)+\tfrac12\log t\log(1-t)$, y aunque
$\Li$ es continua, $\log t\cdot\log(1-t)$ es una indeterminaci\'on $(-\infty)\cdot 0$. Se doma
por compresi\'on: en $(0,\tfrac12]$ se tiene $-\log(1-t)\le 2t$ (de $\log z\ge1-z^{-1}$),
luego
\[
  0\ \le\ \log t\cdot\log(1-t)=(-\log t)\,(-\log(1-t))\ \le\ 2t\,(-\log t)=2\,(-t\log t)
  \ \longrightarrow\ 0,
\]
porque $t\log t\to0$ (el \lean{Real.continuous\_mul\_log} de Mathlib). As\'i $G\to-\Lrog(y)$;
siendo constante e igual a $G(x)$ cerca de $0$, \emph{vale} $-\Lrog(y)$, que es el
Teorema~\ref{thm:five}. \qed

As\'i que la respuesta a Q2 es s\'i: el generador no necesita maquinaria que las identidades
elementales no usaran ---solo que la cancelaci\'on de nueve-en-cinco cierre de verdad.
\emph{Y ah\'i la pluma y la m\'aquina se separaron.}

\section{Q3: ¿qu\'e caz\'o la m\'aquina que la pluma pasa de puntillas?}\label{sec:caught}

\emph{Cuenta hist\'orica.} ``$G'\equiv0$ por un c\'alculo breve'' es como todo manual despacha
el paso~1. La frase esconde un despeje de denominadores en el que, sobre el papel, uno
conf\'ia. Un asistente de pruebas no conf\'ia; lo ejecuta, y la ejecuci\'on tiene un orden.

La cancelaci\'on se verific\'o num\'ericamente a treinta d\'igitos antes de formalizarla ---y
la formalizaci\'on a\'un caz\'o una sutileza que la vista pasa por alto. Cuando
\lean{field\_simp} limpia los denominadores de $G'(t)$ reordena en silencio el producto
$x\cdot y$ a $y\cdot x$, y luego necesita la hip\'otesis $1-y x\neq0$ ---en ese orden exacto
de t\'erminos--- presente en el contexto para descargar el inverso que produce. Aportada un
instante demasiado tarde, el goal no cierra, y ninguna cantidad de \lean{ring} lo repara;
aportada en el orden equivocado ($1-xy$ en vez de $1-yx$), lo mismo. El arreglo es una
l\'inea, pero hallarlo cost\'o m\'as que el resto de la prueba: el hecho de no-anulaci\'on
hubo de enunciarse para ambos \'ordenes y colocarse antes de la simplificaci\'on, no
despu\'es.

\begin{recipebox}{En qu\'e se gana el sueldo un asistente de pruebas}
Sobre el papel, ``despeja denominadores y el resto se cancela'' es una sola frase.
Formalmente es una reescritura ordenada: \lean{field\_simp} puede reordenar un producto
($x y \rightsquigarrow y x$) y luego exigir el hecho de no-anulaci\'on que casa
($1-yx\neq0$) \emph{en ese orden, ya en el contexto}. El diagn\'ostico, archivado para el
pr\'oximo desarrollo de peso~2: si \lean{field\_simp} deja un $(1-a)^{-1}$ sin limpiar, la
hip\'otesis que falta es $(1-a)\neq0$ en el orden de t\'erminos que \lean{field\_simp}
produjo ---aporta tanto $1-xy\neq0$ como su conmutado $1-yx\neq0$.
\end{recipebox}

Esto es toda la matem\'atica nueva, y es honestamente peque\~na: no un teorema sino una
disciplina. Es tambi\'en justo la clase de paso que justifica verificar por m\'aquina una
identidad cuya prueba informal todos ya creen. \emph{Con el generador certificado, ¿qu\'e nos
regala ahora el camino largo del paper compa\~nero?}

\section{Q4: con el pent\'agono certificado, ¿qu\'e cae gratis?}\label{sec:corollaries}

\emph{Cuenta hist\'orica.} Las evaluaciones de la raz\'on \'aurea son de Landen (1780) y
Rogers (1907~\cite{rogers}); su papel moderno como cargas centrales conformes pertenece al
ansatz de Bethe termodin\'amico de los a\~nos 90~\cite{zamolodchikov,nahm,kirillov}. Toda
ruta cl\'asica a ellas pasa por la relaci\'on de cinco t\'erminos; el paper compa\~nero tom\'o
un desv\'io, y ahora la carretera principal est\'a abierta.

La relaci\'on de cinco t\'erminos es una m\'aquina que convierte una elecci\'on de $(x,y)$ en
una relaci\'on cerrada entre cinco valores de $\Lrog$. Alim\'entela con los puntos justos y se
imprime el cat\'alogo del tema (Tabla~\ref{tab:instances}). Dos instancias son especialmente
limpias porque sus argumentos \emph{coinciden}. En $(x,y)=(\tfrac12,\tfrac12)$ los dos
argumentos fraccionarios son ambos $\tfrac13$, dando
$2\Lrog(\tfrac12)=\Lrog(\tfrac14)+2\Lrog(\tfrac13)$. En $(x,y)=(1/\varphi,1/\varphi)$ los tres
argumentos compuestos colapsan \emph{todos} a $1/\varphi^2$ (porque $1-1/\varphi^2=1/\varphi$
y $1-1/\varphi=1/\varphi^2$), de modo que la relaci\'on de cinco t\'erminos dice
$2\Lrog(1/\varphi)=3\Lrog(1/\varphi^2)$ ---y con los valores de la escalera
$\Lrog(1/\varphi)=\pi^2/10$, $\Lrog(1/\varphi^2)=\pi^2/15$, esta es la identidad exacta
$\pi^2/5=\pi^2/5$. La escalera \'aurea que el compa\~nero alcanz\'o por una puerta lateral es,
desde aqu\'i, una sola sustituci\'on.

\begin{table}[h]
\centering\small
\rowcolors{2}{ggband!35}{white}
\begin{tabular}{@{}llll@{}}
\toprule
\rowcolor{ggteal!18} $(x,y)$ & argumentos & la relaci\'on dice & valor exacto \\
\midrule
$(\tfrac12,\tfrac12)$ & $\tfrac14,\tfrac13,\tfrac13$
  & $2\Lrog(\tfrac12)=\Lrog(\tfrac14)+2\Lrog(\tfrac13)$
  & $\Lrog(\tfrac14)+2\Lrog(\tfrac13)=\tfrac{\pi^2}{6}$ \\
$(\tfrac1\varphi,\tfrac1\varphi)$ & $\tfrac1{\varphi^2},\tfrac1{\varphi^2},\tfrac1{\varphi^2}$
  & $2\Lrog(\tfrac1\varphi)=3\Lrog(\tfrac1{\varphi^2})$
  & $\tfrac{\pi^2}{5}=\tfrac{\pi^2}{5}$ \\
\bottomrule
\end{tabular}
\caption{El generador en acci\'on: dos elecciones de $(x,y)$ cuyos argumentos compuestos
coinciden, y la identidad exacta que cada una imprime. Ambas filas reverificadas a treinta
d\'igitos; la primera es comprobable a mano desde $\Lrog(\tfrac12)=\pi^2/12$, la segunda
desde los valores de la escalera $\Lrog(1/\varphi)=\pi^2/10$, $\Lrog(1/\varphi^2)=\pi^2/15$.
La reflexi\'on de Euler $\Lrog(x)+\Lrog(1-x)=\pi^2/6$ \emph{no} es una instancia de la
relaci\'on de cinco t\'erminos sino su compa\~nera ---la simetrizaci\'on que limpia la forma
de Rogers--- y se prueba aparte.}
\label{tab:instances}
\end{table}

\begin{recipebox}{La receta: leer una evaluaci\'on en $\pi^2$ en tres l\'ineas}
Elige cualquier $x,y\in(0,1)$ en los que los cuatro argumentos sean amables. (1) Forma
$xy,\ \tfrac{x(1-y)}{1-xy},\ \tfrac{y(1-x)}{1-xy}$. (2) Aplica
$\Lrog(x)+\Lrog(y)=$ su suma, con garant\'ia verificada por m\'aquina. (3) Sustituye los
valores de $\Lrog$ que conozcas.\\[3pt]
\emph{Trabajado:} toma $x=y=\tfrac12$. Entonces $xy=\tfrac14$ y ambas fracciones son
$\tfrac13$, luego $2\,\Lrog(\tfrac12)=\Lrog(\tfrac14)+2\,\Lrog(\tfrac13)$. Como
$\Lrog(\tfrac12)=\Li(\tfrac12)+\tfrac12\log^2\!2$ y $\Li(\tfrac12)=\tfrac{\pi^2}{12}-
\tfrac12\log^2\!2$, los t\'erminos logar\'itmicos se cancelan y
$\Lrog(\tfrac12)=\tfrac{\pi^2}{12}$; luego $\Lrog(\tfrac14)+2\,\Lrog(\tfrac13)=
\tfrac{\pi^2}{6}$ ---una forma cerrada para tres valores del dilogaritmo, certificada.
\end{recipebox}

As\'i que el hilo que el paper compa\~nero dej\'o colgando queda cerrado. Subi\'o la escalera
\'aurea sin Abel y se pregunt\'o si la puerta lateral era una curiosidad o una necesidad; la
respuesta es que era un atajo, y la escalinata que rode\'o est\'a ahora verificada por
m\'aquina. Las cuatro preguntas terminan donde empezaron ---en una suma de valores del
dilogaritmo igual a $\pi^2/6$--- pero la identidad detr\'as ya no es una sombra en la pared:
es un teorema en la biblioteca, con todo el tema de peso~2 colgando de \'el.

\section{Los bloques de construcci\'on}\label{sec:blocks}

El resultado central descansa sobre tres piezas reutilizables, cada una \lean{sorry}-libre
con los tres axiomas est\'andar (Tabla~\ref{tab:inventory}).

\begin{lemma}[\lean{Dilog.hasDerivAt\_rogersL}]
En $(0,1)$, $\displaystyle\Lrog'(x)=-\tfrac12\Big(\tfrac{\log(1-x)}{x}+\tfrac{\log
x}{1-x}\Big)$.
\end{lemma}
El mismo esquema de derivaci\'on que el compa\~nero us\'o para la reflexi\'on y Landen:
combina $\Li'(x)=-\log(1-x)/x$ con la regla del producto para $\tfrac12\log x\log(1-x)$.

\begin{lemma}[\lean{Dilog.tendsto\_rogersL\_nhdsGT\_zero}]
$\Lrog(t)\to0$ cuando $t\to0^+$.
\end{lemma}
La parte $\Li$ es continua en $[-1,1]$; el producto de logaritmos se comprime a $0$ por
$2(-t\log t)$. Es el \'unico lugar donde la prueba deja el \'algebra pura por el an\'alisis.

\begin{lemma}[\lean{Dilog.continuousAt\_rogersL}]
$\Lrog$ es continua en todo $x\in(0,1)$.
\end{lemma}
Necesario para el \'unico t\'ermino de frontera que sobrevive,
$\Lrog\big(\tfrac{y(1-t)}{1-ty}\big)\to\Lrog(y)$.

\begin{table}[h]
\centering\small
\rowcolors{2}{ggband!35}{white}
\begin{tabular}{@{}lll@{}}
\toprule
\rowcolor{ggteal!18} Nombre Lean & enunciado & d\'onde \\
\midrule
\lean{rogersL\_fiveterm} & la relaci\'on de cinco t\'erminos (Teorema~\ref{thm:five}) & \S\ref{sec:proof} \\
\lean{fiveterm\_hasDerivAt\_zero} & $G'\equiv0$: nueve logs colapsan a cinco & \S\ref{sec:proof}--\ref{sec:caught} \\
\lean{hasDerivAt\_rogersL} & $\Lrog'(x)=-\tfrac12(\tfrac{\log(1-x)}{x}+\tfrac{\log x}{1-x})$ & \S\ref{sec:blocks} \\
\lean{tendsto\_rogersL\_nhdsGT\_zero} & $\Lrog(t)\to0$ cuando $t\to0^+$ & \S\ref{sec:blocks} \\
\lean{continuousAt\_rogersL} & $\Lrog$ continua en $(0,1)$ & \S\ref{sec:blocks} \\
\bottomrule
\end{tabular}
\caption{Los resultados que cargan el peso, todos \lean{sorry}-libres en
\lean{Dilog/FiveTerm.lean} del repositorio, sobre el dilogaritmo de \lean{Dilog/Basic.lean}
del paper compa\~nero. Axiomas de cada entrada: \lean{propext}, \lean{Classical.choice},
\lean{Quot.sound}, y nada m\'as.}
\label{tab:inventory}
\end{table}

\section{El l\'imite de la contribuci\'on}\label{sec:boundary}

Todo lo que aqu\'i es matem\'atica es cl\'asico. La relaci\'on de cinco t\'erminos es de
Spence y Abel; la funci\'on $\Lrog$ de Rogers es de Rogers (1907); los enunciados de
universalidad son de Wojtkowiak y de~Jeu; la prueba por derivaci\'on es est\'andar. No afirmo
ninguna identidad nueva ni ninguna desigualdad nueva. Las contribuciones son tres: la
\emph{formalizaci\'on} (hasta donde s\'e, la primera relaci\'on de cinco t\'erminos
verificada por m\'aquina, sobre el primer dilogaritmo verificado por m\'aquina, en cualquier
asistente de pruebas); la \emph{infraestructura reutilizable} ---la derivada de $\Lrog$ de
Rogers, la continuidad por la derecha en $0$ v\'ia la compresi\'on $-t\log t$, y la
continuidad interior--- que es exactamente lo que cualquier desarrollo de peso~2 posterior
(escaleras, el grupo de Bloch, el cinco t\'erminos de Bloch--Wigner) va a necesitar; y el
\emph{registro de ingenier\'ia de pruebas} de la Secci\'on~\ref{sec:caught}, donde la
cancelaci\'on de nueve-en-cinco es la clase de paso en que una verificaci\'on formal se gana
el sueldo.

\emph{Sobre la metodolog\'ia.} La formalizaci\'on se hizo con un asistente de IA (Claude,
Anthropic), bajo la disciplina del paper compa\~nero: la cancelaci\'on central verificada
num\'ericamente a treinta d\'igitos antes de formalizar, y cada teorema comprobado con
\lean{\#print axioms} para depender solo de \lean{propext}, \lean{Classical.choice} y
\lean{Quot.sound}. La correcci\'on no descansa en el asistente: el n\'ucleo de Lean verifica
cada prueba. La matem\'atica y las afirmaciones son responsabilidad del autor.

\section{Preguntas abiertas}\label{sec:questions}

El pent\'agono es ahora un teorema en la biblioteca. De pie sobre \'el, seis preguntas que
merecen el pr\'oximo esfuerzo ---la \'ultima es la que yo me quedar\'ia.

\begin{enumerate}
\item \emph{Las escaleras como familia generada.} Cada fila de la Tabla~\ref{tab:instances}
es la relaci\'on de cinco t\'erminos evaluada en un punto; mecanizar la sustituci\'on para que
cada escalera \emph{se imprima} desde el \'unico teorema convertir\'ia un cat\'alogo en un
corolario. ¿Cu\'anta de la tabla cl\'asica de valores del dilogaritmo es una \'orbita finita
de la relaci\'on de cinco t\'erminos bajo sustituci\'on $\mathbb{Q}$-racional ---y d\'onde se
detiene esa \'orbita?
\item \emph{La secuela de Bloch--Wigner.} La misma relaci\'on vale para la funci\'on
univaluada de Bloch--Wigner en $\mathbb{C}$, donde sostiene el volumen hiperb\'olico y la
congruencia por tijeras~\cite{suslin,neumannzagier,goncharov}. ¿Cu\'al es el menor relato
fiel del grupo de Bloch en Lean en el que $\mathrm{Vol}$ de una $3$-variedad hiperb\'olica se
vuelve un n\'umero computable y certificado?
\item \emph{El l\'imite cu\'antico.} El dilogaritmo cu\'antico de Faddeev--Kashaev satisface
una identidad del pent\'agono~\cite{faddeevkashaev} cuyo l\'imite $\hbar\to0$ es la relaci\'on
probada aqu\'i. ¿Puede un asistente de pruebas hacer de ese l\'imite un teorema en vez de una
heur\'istica ---y, al hacerlo, expondr\'ia, como lo hizo la prueba cl\'asica, un paso que la
literatura f\'isica pasa de puntillas?
\item \emph{La universalidad, formalizada.} Wojtkowiak y de~Jeu reducen \emph{toda} ecuaci\'on
funcional racional de $\Li$ a la relaci\'on de cinco t\'erminos~\cite{wojtkowiak,dejeu}. Esa
reducci\'on es un algoritmo. ¿Qu\'e har\'ia falta para correrlo dentro de Lean, de modo que
una identidad de dilogaritmo cualquiera se \emph{decida} contra el generador?
\item \emph{Lo que ve la m\'aquina y la p\'agina esconde.} La Secci\'on~\ref{sec:caught}
hall\'o un lugar donde la prueba formal fue estrictamente m\'as afilada que la informal ---una
reescritura ordenada en la que la pluma conf\'ia. En un tema de dos siglos, reprobado una
docena de veces, ¿cu\'antos de esos pasos silenciosos quedan, y es ``la prueba que un humano
cree'' alguna vez exactamente la prueba que una m\'aquina aceptar\'a?
\item \emph{Una para guardar.} La relaci\'on de cinco t\'erminos dice que cinco valores del
dilogaritmo, le\'idos en torno a un pent\'agono, suman nada nuevo. Es la relaci\'on que define
el grupo de Bloch, la mutaci\'on de cl\'uster, el dilogaritmo cu\'antico ---una identidad con
las m\'ascaras de la geometr\'ia, el \'algebra y la f\'isica. Ahora que est\'a verificada por
m\'aquina, la pregunta no es si es cierta sino por qu\'e est\'a \emph{en todas partes}: ¿qu\'e
tiene un pent\'agono para que tantas partes de la matem\'atica sean, por debajo, los mismos
cinco pasos que te devuelven a casa?
\end{enumerate}

\section*{Reproducir}
\begin{verbatim}
git clone https://github.com/karlesmarin/godsil-gutman-lean
cd godsil-gutman-lean && lake exe cache get
lake build Dilog.FiveTerm
\end{verbatim}
Lean \lean{v4.30.0-rc2}, Mathlib pin \lean{701fb6e9}. La relaci\'on es
\lean{Dilog.rogersL\_fiveterm} en \lean{Dilog/FiveTerm.lean}; sus soportes est\'an en la
Tabla~\ref{tab:inventory}. Cada uno se comprueba con \lean{\#print axioms} para depender solo
de \lean{propext}, \lean{Classical.choice}, \lean{Quot.sound}.

\begin{thebibliography}{99}\small
\bibitem{abel} N.~H.~Abel, \emph{\OE uvres compl\`etes} (L.~Sylow, S.~Lie, eds.), vol.~II,
Gr{\o}ndahl, Christiania, 1881; la relaci\'on de cinco t\'erminos est\'a en la nota sobre el
dilogaritmo.
\bibitem{rogers} L.~J.~Rogers, On function sum theorems connected with the series
$\sum x^n/n^2$, \emph{Proc.\ London Math.\ Soc.} \textbf{4} (1907) 169--189.
\bibitem{lewin} L.~Lewin, \emph{Polylogarithms and Associated Functions}, North-Holland,
1981.
\bibitem{zagier} D.~Zagier, The dilogarithm function, en \emph{Frontiers in Number Theory,
Physics, and Geometry II}, Springer, 2007, 3--65.
\bibitem{wojtkowiak} Z.~Wojtkowiak, The basic structure of polylogarithmic functional
equations, en \emph{Structural Properties of Polylogarithms} (L.~Lewin, ed.), Math.\
Surveys Monogr.\ \textbf{37}, AMS, 1991, 205--231.
\bibitem{dejeu} R.~de~Jeu, Towards regulator formulae for the $K$-theory of curves and
number fields, \emph{Compositio Math.} \textbf{124} (2000) 137--194.
\bibitem{bloch} S.~Bloch, \emph{Higher Regulators, Algebraic $K$-theory, and Zeta Functions
of Elliptic Curves}, CRM Monogr.\ Ser.\ \textbf{11}, AMS, 2000 (lecciones de Irvine, 1978).
\bibitem{suslin} A.~A.~Suslin, $K_3$ of a field, and the Bloch group, \emph{Proc.\ Steklov
Inst.\ Math.} \textbf{183} (1991) 217--239.
\bibitem{neumannzagier} W.~D.~Neumann, D.~Zagier, Volumes of hyperbolic three-manifolds,
\emph{Topology} \textbf{24} (1985) 307--332.
\bibitem{goncharov} A.~B.~Goncharov, Geometry of configurations, polylogarithms, and
motivic cohomology, \emph{Adv.\ Math.} \textbf{114} (1995) 197--318.
\bibitem{nakanishi} T.~Nakanishi, Dilogarithm identities for conformal field theories and
cluster algebras, en \emph{Symmetries, Integrable Systems and Representations}, Springer
Proc.\ Math.\ Stat.\ \textbf{40} (2013) 407--444.
\bibitem{faddeevkashaev} L.~D.~Faddeev, R.~M.~Kashaev, Quantum dilogarithm as a kernel of
the quantum modular double, \emph{Modern Phys.\ Lett.\ A} \textbf{9} (1994) 427--434.
\bibitem{kirillov} A.~N.~Kirillov, Dilogarithm identities, \emph{Prog.\ Theor.\ Phys.\
Suppl.} \textbf{118} (1995) 61--142.
\bibitem{zamolodchikov} Al.~B.~Zamolodchikov, Thermodynamic Bethe ansatz in relativistic
models, \emph{Nucl.\ Phys.\ B} \textbf{342} (1990) 695--720.
\bibitem{nahm} W.~Nahm, Conformal field theory and torsion elements of the Bloch group, en
\emph{Frontiers in Number Theory, Physics, and Geometry II}, Springer, 2007, 67--132.
\bibitem{clock} C.~Mar\'in, \emph{The Clock That Never Ticks: a machine-checked path from
the dilogarithm to quantum speed limits in Lean~4}, paper compa\~nero, mismo repositorio,
2026.
\bibitem{mathlib} The mathlib Community, The Lean mathematical library, en \emph{CPP 2020},
ACM, 2020, 367--381.
\end{thebibliography}

\end{document}
